% !TEX root = ../aa_SoM_Chapters.tex

%% HARRIS: Chapter 10

% ö = \%c3\%b6
% ä = \%C3\%A4

\xdef\sectiontitle{\IIsectiontitle} %aiheuttama

%%%%%%%%%%%%%%%
\begin{frame}[label=2_10_a]%[label=2_9_TOC1]
\frametitle{\texorpdfstring{\culine[\sectiontitleulinecolor]{\sectiontitle}}{\sectiontitle} }

%\vspace{0.5cm}
%\begin{itemize}[leftmargin=0.5cm,itemsep=3mm]
%    \item[2.1] \hyperlink{2_1_a}{\IIaSubsectiontitle}
%    \item[2.2] \hyperlink{2_2_a}{\IIbSubsectiontitle}	
%    \item[2.3] \hyperlink{2_3_a}{\IIcSubsectiontitle}	
%    \item[2.4] \hyperlink{2_4_a}{\IIdSubsectiontitle}
%    \item[2.5] \hyperlink{2_5_a}{\IIeSubsectiontitle}
%    \item[2.6] \hyperlink{2_6_a}{\IIfSubsectiontitle}
%    \item[2.7] \hyperlink{2_7_a}{\IIgSubsectiontitle}
%    \item[2.8] \hyperlink{2_8_a}{\IIhSubsectiontitle}
%    \item[2.9] \hyperlink{2_9_a}{\CDAlert[blue]{\IIiSubsectiontitle}}
%\end{itemize}

\vspace{0.2cm}
\begin{itemize}[leftmargin=0.5cm,itemsep=3mm]
    \item[2.1] \hyperlink{2_1_a}{\IIaaSubsectiontitle}
    \item[2.2] \hyperlink{2_2_a}{\IIaSubsectiontitle}
    \item[2.3] \hyperlink{2_3_a}{\IIbSubsectiontitle}	
    \item[2.4] \hyperlink{2_4_a}{\IIcSubsectiontitle}	
    \item[2.5] \hyperlink{2_5_a}{\IIdSubsectiontitle}
    \item[2.6] \hyperlink{2_6_a}{\IIeSubsectiontitle}
    \item[2.7] \hyperlink{2_7_a}{\IIfSubsectiontitle}
    \item[2.8] \hyperlink{2_8_a}{\IIgSubsectiontitle}
    \item[2.9] \hyperlink{2_9_a}{\IIhSubsectiontitle}
    \item[2.10] \hyperlink{2_10_a}{\CDAlert[blue]{\IIiSubsectiontitle}}
\end{itemize}

\endofslide \end{frame}
%%%%%%%%%%%%%%%

\xdef\sectiontitle{\IIiSubsectiontitle} 

%%%%%%%%%%%%%%%
\begin{frame}
\frametitle{\texorpdfstring{\culine[\sectiontitleulinecolor]{\sectiontitle}}{\sectiontitle} }

%On fullerenes and Carbon nanotubes (CNT)
\begin{itemize}
	\item Progress in material science has come up with a multitude of 3D\appear{, 2D} 
\end{itemize}
%
%\vspace{1.6mm}
\begin{minipage}{6.5cm}
\begin{itemize}%[itemsep=1.6mm]
	\item[] and even 1D carbon nanostructures \appear{such as fullerenes}\appear{, \CDAlert[red]{graphene}}\appear{, and \CDAlert[blue]{carbon nanotubes (CNTs)} }
	
	\item Performance of CNTs can be improved by doping with functional atoms/molecules \appear{(\CDAlert[blue]{optoelectronics})}

\item Fullerenes have been also used as molecular ball bearings\appear{, microscopic strainers} \appear{and as nanoscale size standards}

\item CNT's and graphene have high mechanical strength
\implies{ Sports equipment and protective clothing }
	\end{itemize}
\end{minipage}

\xdef\loc{south east}
\begin{tikzpicture}[remember picture,overlay] % anchor=south
    \node (A) [yshift=-0.25*\figYshift,xshift=\figXshift, anchor=\loc] at (current page.\loc) {\includegraphics[page=1,width=7cm]{Figures_booklet/SOM_10_Bonding_figures/SOM_Bonding_Figure_57.png}}; %scale=\figscale. % 8cm max height
\end{tikzpicture}

%% 58 and 61 figs too
\endofslide 
%\endofsection
\end{frame}
%%%%%%%%%%%%%%%

%%%%%%%%%%%%%%%
\begin{frame}
\frametitle{\texorpdfstring{\culine[\sectiontitleulinecolor]{\sectiontitle}}{\sectiontitle} }

\begin{itemize}
	\item A special group of achiral single-walled CNTs are \Alert[red]{metallic}\appear{, but all the rest of the CNTs are either \Alert[blue]{small or moderate band gap semiconductors}}\appear{, thus they give a high promise to be utilized in nanoelectronics} \appear{(\CDAlert[fuchsia]{FETs} and other components)} \appear{and e.g. as nanoantenna to detect light waves (analogy to radio antenna)}
\end{itemize}

\xdef\loc{south west}
\begin{tikzpicture}[remember picture,overlay] % anchor=south
    \node (A) [yshift=-0.25*\figYshift,xshift=-\figXshift, anchor=\loc] at (current page.\loc) {\includegraphics[page=1,width=7cm]{Figures_booklet/SOM_10_Bonding_figures/SOM_Bonding_Figure_59.png}}; %scale=\figscale. % 8cm max height
\end{tikzpicture}

\xdef\loc{south east}
\begin{tikzpicture}[remember picture,overlay] % anchor=south
    \node (A) [yshift=-0.25*\figYshift,xshift=\figXshift, anchor=\loc] at (current page.\loc) {\includegraphics[page=1,width=7cm]{Figures_booklet/SOM_10_Bonding_figures/SOM_Bonding_Figure_60.png}}; %scale=\figscale. % 8cm max height
\end{tikzpicture}

\endofslide 
%\endofsection
\end{frame}
%%%%%%%%%%%%%%%



